\section{Conditions, parity, and telemetry}
\label{sec:conditions}

\paragraph{The four rungs.}
All four conditions are reasonably developed on training data, frozen together,
and evaluated in the same sealed event. Table~\ref{tab:conditions} states what
each changes.

\begin{table}[htbp]
\centering
\caption{The four frozen conditions as specified. The enforcement column states
the design intent; what the governed condition did at run time is measured later
in this section.}
\label{tab:conditions}
\begin{tabular}{p{2.3cm}p{3.1cm}p{3.6cm}p{3.6cm}}
\toprule
Condition & Knowledge & Representation & Enforcement \\
\midrule
C1 Raw SQL & Public schema & Raw schema & None \\
C2 HKB reference & Public schema and HKB & Searchable raw HKB & Optional \\
C3 Model reference & Public schema and HKB & Searchable exported semantic model & Optional \\
C4 Governed & Public schema and HKB & Semantic model & Production harness enforces semantic compilation and validation \\
\bottomrule
\end{tabular}
\end{table}

C2 and C3 receive comparable programmatic search rather than being made to ingest
a roughly 84K-token export in one prompt; the difference between them is the
representation, not the retrieval budget. C1 through C3 share one pinned
direct-SQL model and harness and match question wording, database, execution
access, budgets, retries, and scorer. C2 is a substantive condition and not a
bridge: business semantics supplied as searchable context may account for a
material part of the total C4 minus C1 gap, and evidence about semantic context
informs C2 minus C1 only.

C4 uses the production-default configuration, because production fidelity
outranks forcing model parity for the primary product endpoint. The exact C4
models could not be observed and invoked independently: model version, retry
count, and validation-attempt count are unavailable for every held-out C4
attempt. Provider, observable identifiers, and settings are recorded where the
product exposes them, and the interpretation is therefore the conservative one
stated in advance. C4 minus C3 is a production-system comparison and is not the
causal effect of enforcement. All conditions receive real
development rather than leaving C1 through C3 as straw comparators; per condition
we record prompt revisions, training evaluations, major interventions,
approximate effort, and the actor or automation that performed the work. Exact
equal person-hours are not required.

\paragraph{Parity policy.}
Pre-treatment resources are matched where matching is meaningful: database
snapshot, credentials and permissions, public information, task specification,
model identity when possible, reasoning and token ceilings, and retry ceilings.
Post-treatment behavior is not artificially equalized. Actual tool calls, tokens
consumed, latency, database-query count, retries, and cost are measured as
outcomes, because a governed system that spends more to be right and a direct
agent that spends less to be wrong differ in a way an equalized budget would
hide.

\paragraph{Telemetry contract.}
Every attempt produces an immutable generation envelope containing no correctness
information. It records run and attempt identifiers, condition, repetition,
partition, and timestamps; provider, model, and version; token counts and their
source; cost and its source; wall-clock latency, tool calls, database queries,
retries, and validation attempts; the generation outcome
(\texttt{answered}, \texttt{refused}, or \texttt{errored}), terminal failure
origin and class, the generated SQL or governed query, semantic objects, and
fixed compiler, validation, and execution statuses; and a relative path, schema
version, and SHA-256 for the ordered raw trace, or an explicit degraded reason
when trace capture is unavailable.

Unobserved counts are null and must be named in an explicit unavailability list;
they are never encoded as zero. Token and cost values separately declare whether
they are provider-reported, derived, or unavailable. Complete traces must
reconcile token, tool-call, database-query, retry, and validation totals to the
attempt envelope, and truncated traces carry an explicit degraded reason. Scoring
produces a \emph{separate} record carrying \texttt{correct},
\texttt{wrong\_answer}, or \texttt{refused\_or\_error}, hash-bound to both the
generation file and every attempt record; it cannot mutate or copy the prompt,
query, result, or telemetry. Raw traces and generated SQL live only under ignored
run roots, and committed artifacts contain hashes and permitted summaries.

Two consequences are worth stating explicitly. The three-state attempt outcome is
distinct from the three-state generation outcome, so a confident wrong answer and
a safe refusal are never collapsed even though both miss the benchmark. And for
composite production workflows, an aggregate model label must not be presented as
exact model parity when internal routing is opaque; where the product exposes
per-stage model and usage, raw capture preserves them.

\paragraph{Capture verification gate.}
Before any 231-question baseline or expensive experiment, each implemented
condition had to produce one public unscored smoke attempt whose envelope
validates, whose trace hashes and redaction verify (or whose unexposed fields
carry an explicit degraded reason), whose system database-query counts exclude
scorer and gold execution, and whose evaluated-system retries are distinguishable
from benchmark-infrastructure reruns. Coverage was recorded per field and
condition. The gate passed for all four conditions, and the field-level coverage
it recorded is what makes the unavailability statements in
Section~\ref{sec:results} auditable rather than rhetorical: latency coverage is
complete in every condition, token, tool-call, and database-query coverage is
267 of 267 attempts for C1 through C3 and 265 of 267 for C4, and
provider-reported dollar cost is 267 of 267 for C1 through C3 and 0 of 267 for
C4.

\paragraph{What the governed condition turned out to be.}
C4 was preregistered as the condition in which the production harness enforces
semantic compilation and validation. Measured on the frozen development
baseline, it is also an agent authoring SQL. All 135 captured governed semantic
queries carry \texttt{rewriteSql: true} and \texttt{aiGenerated: true} with
hand-authored SQL in \texttt{userEditedSQL}; \texttt{generated\_sql} is null on
every attempt; zero queries declare a join path; and \texttt{join\_via\_map} is
empty throughout. The choice of path was the product's own. The harness posts
only a model identifier, the bare question, and a branch identifier, and exposes
no mode flag; a repository-wide search establishes that no harness code ever
writes \texttt{rewriteSql}.

The reason the path existed is in the deployed model rather than in the agent.
Because the conservative public-only compiler deferred 46.9\% of HKB definitions
across an unresolved grain (Section~\ref{sec:development}), the deployed topics
publish no joins and no measures, so for any cross-table or aggregate question no
compiled path was available to take. The semantic layer still did work as a field
vocabulary that the product resolved at rewrite time: 109 of 135 attempts
reference at least one compiled dimension, and 39 expand an HKB-backed derived
definition into the executed query. It did not compose the query. Two
consequences follow and are carried through the rest of the paper. C4 minus C3
does not isolate semantic-layer query composition, and no achievable model parity
would make it do so, because neither arm resolves a join or compiles a measure.
And C4's structural aggregates are computed from \texttt{userEditedSQL}, so they
describe agent-written queries in all four conditions. The full measurement,
including every prior claim it retracts, is in the query-path disclosure listed
in Section~\ref{sec:materials}.
